\section{Literature Review and Background Knowledge}
\label{sec:litreview}

\subsection{Overview}

This section establishes the background needed to understand the project. It
begins with the function of a Security Operations Centre and the role of the
analyst, then examines the technology that underpins it: Security Information
and Event Management (SIEM) platforms, log management, and endpoint telemetry.
It compares open-source and commercial SIEM options, outlines the main
approaches to threat detection and the MITRE ATT\&CK framework, and considers
the use of virtualised home labs for developing practical skills. The section
closes with a review of related work and a statement of how this project
differs from it.

\subsection{Security Operations Centres}

A Security Operations Centre (SOC) is the organisational function responsible
for continuously monitoring an environment, detecting malicious or anomalous
activity, and coordinating the response to security incidents. Rather than a
single technology, a SOC is best understood as the combination of people,
process, and tooling working together to maintain an organisation's security
posture.

SOC analysts are commonly organised into tiers. Tier~1 analysts perform the
initial triage of incoming alerts, deciding quickly whether an alert
represents genuine malicious activity (a true positive) or harmless behaviour
that merely resembles it (a false positive). Alerts that cannot be dismissed
are escalated to Tier~2 analysts for deeper investigation, and the most
serious or complex cases reach Tier~3 staff who carry out incident response
and proactive threat hunting. This triage-first model means that the largest
share of day-to-day SOC work consists of rapidly reviewing, classifying, and
documenting alerts.

The structured handling of incidents that a SOC supports is described in
NIST's incident-handling guidance, which frames incident response as a
lifecycle of preparation; detection and analysis; containment, eradication and
recovery; and post-incident activity~\cite{nist80061}. Detection and analysis
--- the stage at which raw data becomes an actionable alert --- is the stage
this project is most concerned with, and it depends almost entirely on the
SOC's central tooling: the SIEM.

\subsection{Security Information and Event Management}

A SIEM platform is the technological core of a SOC. The term combines two
earlier categories: Security Information Management (SIM), concerned with the
long-term collection, storage, and analysis of log data, and Security Event
Management (SEM), concerned with the real-time monitoring and correlation of
events. A modern SIEM unifies both.

In practice, a SIEM performs several core functions. It \textit{collects} log
and event data from many different sources across an estate. It
\textit{normalises} that data, parsing records written in many different
formats into a consistent structure so they can be compared and searched. It
\textit{correlates} events, applying rules that recognise patterns spanning
multiple log sources --- for example, linking a series of failed logins
followed by a success into a single suspicious sequence. It \textit{alerts}
analysts when activity matches a rule of concern, and it provides
\textit{dashboards and search} so that analysts can investigate. Finally, it
\textit{retains} data for the periods required by forensic investigation and
regulatory compliance.

SIEM platforms have evolved considerably. Early systems were primarily log
aggregators; later generations added correlation and alerting; current
platforms increasingly incorporate user and entity behaviour analytics,
threat-intelligence enrichment, and automated response. Throughout this
evolution, the underlying premise has remained constant: an organisation
cannot detect what it does not collect and analyse, and the SIEM is the
instrument through which that collection and analysis happens.

\subsection{Log Management and Endpoint Telemetry}

Because detection ultimately rests on log data, the quality of that data is
decisive. NIST's guidance on computer security log management emphasises that
logs are essential not only for detecting attacks but also for understanding
them after the fact, and that effective log management requires deliberate
planning of what to collect and how to handle it~\cite{nist80092}. National
guidance makes a similar point: logging must be designed with security
outcomes in mind rather than left to operating-system defaults~\cite{ncsc}.

This distinction --- between default logging and security-grade logging --- is
central to the present project. On Windows, the standard Event Log captures
useful records such as logon events and service changes, but it does not by
itself provide the fine-grained visibility that behavioural detection
requires. \textbf{Sysmon} (System Monitor), a free tool from Microsoft's
Sysinternals suite, addresses this by logging detailed information about
process creation, network connections, image loads, file-creation timestamps,
and related activity, written to a dedicated event channel~\cite{sysmon}. This
additional telemetry is what allows a SIEM to reason about \textit{behaviour}
rather than merely about discrete system messages.

Linux provides an equivalent capability through the Linux Audit framework and
its userspace daemon, \texttt{auditd}. By defining audit rules, an
administrator can record security-relevant system activity --- notably the
execution of commands through the \texttt{execve} system call --- at the
kernel level. As with Sysmon on Windows, this elevates the endpoint from
producing general operational logs to producing detection-grade telemetry.

The recurring principle across both platforms is that endpoint visibility is
something that must be deliberately engineered. A SIEM connected only to
default logs will see far less than one fed by purpose-configured endpoint
telemetry.

\subsection{Open-Source SIEM Platforms}

SIEM capability is available both as commercial products and as open-source
software. Commercial platforms such as Splunk, Microsoft Sentinel, and IBM
QRadar are powerful and widely deployed, but they are typically licensed by
data volume, which can make them costly --- and largely inaccessible for
independent learning and experimentation.

Several capable open-source alternatives exist, including the Elastic Stack
(often deployed as Elastic Security), Security Onion, Graylog, and Wazuh. This
project uses \textbf{Wazuh}. Wazuh originated as a fork of
\textbf{OSSEC}, a long-established open-source host-based intrusion detection
system~\cite{ossec}, and has since grown into a full SIEM and extended
detection and response (XDR) platform~\cite{wazuhdocs}. A Wazuh deployment
consists of a central \textit{manager} that analyses incoming data and applies
detection rules, an \textit{indexer} that stores and searches events, a
\textit{dashboard} that presents them, and lightweight \textit{agents}
installed on monitored endpoints. Beyond log analysis and alerting, Wazuh
provides file-integrity monitoring, configuration assessment, vulnerability
detection, and the mapping of alerts to the MITRE ATT\&CK framework.

For the purpose of building practical skills, an open-source platform offers
clear advantages. It can be deployed at no cost, it can be examined and
reconfigured freely, and --- importantly --- the concepts it teaches, such as
agent-based collection, log parsing, rule-based correlation, and alert triage,
transfer directly to commercial platforms. The specific product differs; the
underlying discipline does not.

\subsection{Threat Detection Approaches and the MITRE ATT\&CK Framework}

SIEM platforms detect threats using two broad approaches.
\textbf{Signature- or rule-based} detection matches activity against known
patterns of malicious behaviour; it is precise and explainable, but it can
only catch what has been defined in advance. \textbf{Anomaly-based} detection
flags deviations from an established baseline of normal behaviour; it can in
principle catch novel activity, but it is prone to raising alerts on benign
deviations. NIST's guidance on intrusion detection and prevention systems
describes both approaches and notes that practical deployments generally
combine them, since each compensates for the other's
weakness~\cite{nist80094}.

A persistent operational problem for SOCs is \textit{alert fatigue}: when
detection produces a high volume of low-value or false-positive alerts,
analysts become desensitised, and genuine threats risk being overlooked among
the noise. This makes the \textit{tuning} of detection --- raising true
positives while suppressing false positives --- as important as detection
itself.

To make detection systematic, the security community has increasingly adopted
the \textbf{MITRE ATT\&CK} framework: a structured, continually updated
knowledge base of adversary tactics and techniques drawn from real-world
observations of attacks~\cite{mitreattack}. ATT\&CK gives defenders a common
vocabulary and a way to measure detection \textit{coverage} --- mapping which
adversary techniques an environment can and cannot detect. While the detailed
application of ATT\&CK belongs to the planned follow-on project on detection
engineering, it is introduced here because it is the framework against which
modern detection capability is now routinely assessed.

\subsection{Virtualised Security Labs and the Skills Gap}

The cybersecurity profession is widely recognised to face a shortage of
skilled practitioners, and industry reporting consistently identifies a
mismatch between the volume of security-relevant activity organisations must
handle and the workforce available to handle it~\cite{verizondbir}. For new
entrants specifically, the gap is often less about theoretical knowledge than
about practical, tool-level experience.

Virtualised home labs are a well-established response to this problem. Using a
hypervisor such as VirtualBox or VMware, a learner can build a self-contained
environment of multiple virtual machines on a single physical computer, at no
cost and in complete isolation from any production network. Such a lab can be
broken, rebuilt, and experimented on freely. For defensive security in
particular, a home lab allows a learner to stand up a SIEM, generate security
events safely, and practise the monitoring and triage workflow that a SOC role
demands --- experience that is otherwise difficult to obtain before
employment. Frameworks such as the CIS Critical Security Controls reinforce
the value of this hands-on focus by expressing good security practice in terms
of concrete, implementable actions rather than abstract principles~\cite{ciscontrols}.

\subsection{Related Works}

Building security home labs is an established practice, and several reference
projects exist. The most prominent is \textbf{DetectionLab} by Chris
Long~\cite{detectionlab}, an automated project that provisions a pre-built
Windows domain instrumented with logging and monitoring tooling. DetectionLab
is comprehensive and highly automated; its strength --- one-command
provisioning --- is also a limitation for a learner, as the automation
abstracts away the very build steps from which much of the understanding
comes. A large body of community material, including vendor tutorials and
blog write-ups, similarly documents Wazuh and Elastic-based lab builds, though
such material varies in depth and is rarely presented as a structured,
reproducible whole.

This project is positioned deliberately against that backdrop, and differs in
three respects. First, \textbf{the build is performed and documented manually
and in full}: reproducibility is treated as a primary deliverable, with each
step recorded so that the understanding --- not only the result --- is
transferable. Second, \textbf{the evaluation is explicitly honest about
realism}: rather than presenting a home lab as equivalent to a production SOC,
the report sets out clearly where a single-machine virtualised environment
does and does not represent real operations. Third, \textbf{the project is
designed as the first stage of a two-part progression}: it establishes the
monitoring foundation on which a planned second project, focused on detection
engineering with MITRE ATT\&CK and Atomic Red Team~\cite{atomicredteam}, will
build. The contribution is therefore not a novel tool, but a clear,
honest, and reproducible account of how a credible SOC monitoring environment
is constructed and operated using open-source software.

% NOTE TO AUTHOR: The Related Works subsection currently cites only sources
% that can be verified. To strengthen it for an academic standard, add 2-3
% peer-reviewed papers (e.g. studies evaluating Wazuh or open-source SIEM
% detection performance) found through your own literature search. Do not add
% any reference you have not personally read and verified.
